\chapter{Notación y símbolos}

La siguiente tabla fija la notación principal de esta edición. Cuando una fuente use una letra distinta, se traduce a este sistema antes de continuar.

\begin{longtable}{>{\centering\arraybackslash}p{.10\textwidth}p{.27\textwidth}p{.39\textwidth}p{.09\textwidth}}
\toprule
\textbf{Símbolo} & \textbf{Nombre} & \textbf{Significado} & \textbf{SI}\\
\midrule\endhead
$N$ & esfuerzo axil & Resultante normal a la sección; positivo a tracción & N\\
$V$ o $C_y$ & esfuerzo cortante & Resultante transversal en el plano de análisis & N\\
$M$ o $M_z$ & momento flector & Momento interno asociado a flexión & N m\\
$T$ o $M_x$ & momento torsor & Momento alrededor de la directriz & N m\\
$q(x)$ & carga distribuida & Fuerza por unidad de longitud & N/m\\
$\sigma$ & tensión normal & Fuerza normal local por unidad de área & Pa\\
$\tau$ & tensión tangencial & Fuerza tangencial local por unidad de área & Pa\\
$\varepsilon$ & deformación normal & Cambio relativo de longitud & 1\\
$\gamma$ & deformación angular ingenieril & Cambio angular local & rad\\
$E$ & módulo de Young & Rigidez constitutiva uniaxial & Pa\\
$G$ & módulo de cortadura & Rigidez constitutiva a cortadura & Pa\\
$A$ & área de sección & Área resistente & m$^2$\\
$I$ & segundo momento de área & Propiedad geométrica para flexión & m$^4$\\
$J$ & constante/término torsional & Propiedad geométrica asociada a torsión & m$^4$\\
\bottomrule
\end{longtable}
